\documentclass[11pt]{article}
\usepackage[T1]{fontenc}
\usepackage{amsmath,amssymb}
\usepackage{cite}
\usepackage{booktabs}
\usepackage{graphicx}
\usepackage{xurl}
\usepackage{hyperref}
\hypersetup{
  colorlinks=true,
  linkcolor=blue,
  urlcolor=blue,
  citecolor=blue,
  pdfauthor={Jabulani Monchwe, Martin Chanza, Olorato Motitswane, Mfundo Monchwe}
}
\newcommand{\software}{\textsf{ReliableML}}
\newcommand{\code}[1]{\texttt{\small #1}}
\newcommand{\jabulaniauthor}{
\IEEEauthorblockN{Jabulani Monchwe}
\IEEEauthorblockA{\textit{Applied Business Analytics and Decision Making}\\
\textit{North-West University}\\
\textit{IDEA-Research Club Chairperson}\\
Johannesburg, South Africa\\
\href{mailto:monchwej.25@gmail.com}{monchwej.25@gmail.com}}
}
\newcommand{\martinchanzauthor}{
\IEEEauthorblockN{Martin Chanza}
\IEEEauthorblockA{\textit{Applied Business Analytics and Decision Making}\\
\textit{North-West University}\\
Mafikeng, South Africa\\
\href{mailto:martin.chanza@nwu.ac.za}{martin.chanza@nwu.ac.za}}
}
\newcommand{\oloratoauthor}{
\IEEEauthorblockN{Olorato Motitswane}
\IEEEauthorblockA{\textit{Eduvos Midrand Campus -- IT Faculty Data Science, South Africa}\\
\textit{IDEA-Research Club}\\
\href{mailto:oloratomotitswane@gmail.com}{oloratomotitswane@gmail.com}}
}
\newcommand{\mfundoauthor}{
\IEEEauthorblockN{Mfundo Monchwe}
\IEEEauthorblockA{\textit{Eduvos Midrand Campus -- IT Faculty Data Science, South Africa}\\
\textit{IDEA-Research Club}\\
\href{mailto:diditmfundo@gmail.com}{diditmfundo@gmail.com}}
}
\newcommand{\minenhleauthor}{
\IEEEauthorblockN{Minenhle Mchunu}
\IEEEauthorblockA{\textit{Applied Business Analytics and Decision Making}\\
\textit{Eduvos Midrand Campus -- IT Faculty Data Science, South Africa}\\
\textit{IDEA-Research Club}\\
\href{mailto:Mchunu.minenhle02@gmail.com}{Mchunu.minenhle02@gmail.com}}
}
\newcommand{\projecturl}{\url{https://github.com/Jabualani/ML-Reliabilty}}

\title{\software{} Method Note: Classification and Statistical Evaluation}
\author{\contactauthors}
\date{Version 0.1.1}
\begin{document}
\maketitle
\begin{abstract}
Labeled evaluation separates predictive performance from input-distribution
monitoring. The report combines classification metrics, association tests,
train--test comparisons, and an empirical repeated-cross-validation
bias--variance summary.
\end{abstract}

\section{Classification performance}
\code{evaluate(X\_test, y\_test)} predicts on the supplied rows and compares
predictions with true labels. The classification summary includes a confusion
matrix, accuracy, balanced accuracy, per-class precision, recall, specificity
and F1, macro/support-weighted summaries, and multiclass Matthews correlation
coefficient (MCC). The confusion matrix and class denominators help expose
class imbalance concealed by aggregate accuracy.

\section{Statistical tests}
The report keeps two questions distinct:
\begin{enumerate}
  \item \textbf{Feature--target association:} ANOVA for numeric features,
  joint MANOVA when suitable, and chi-square tests for categorical features
  across target classes.
  \item \textbf{Training versus test:} ANOVA and joint MANOVA for numeric
  features, and chi-square homogeneity tests for categorical features.
\end{enumerate}
Results contain statistics, p-values, degrees of freedom where applicable,
and cautions or an explicit unavailable reason. A rank-deficient MANOVA
residual covariance is reported as unavailable rather than forced into a
singular matrix solution.

\section{Empirical bias--variance}
An optional repeated stratified cross-validation procedure estimates a
finite-sample decomposition of one-hot/Brier prediction error. Probability
outputs are preferred; hard predictions are one-hot encoded as a fallback.
The default five folds and five repeats imply 25 refits, with fewer folds
when the smallest class requires it. The result is not a population
irreducible-noise estimate and can be computationally expensive.

\section{Assumptions and interpretation}
Statistical tests require adequate sample sizes, complete observations where
required, and their usual distributional assumptions. The tests describe
associations and cohort differences; small p-values are not causal evidence.
P-values are sensitive to sample size and multiple testing. Analyze test data
that is representative and held out from model selection. For repeated,
grouped, or temporal observations, ordinary independent-row uncertainty
estimates may be too narrow.

\section{Interface}
\begin{verbatim}
evaluation = model.evaluate(X_test, y_test)
print(evaluation.classification["matthews_correlation_coefficient"])
print(evaluation.classification["confusion_matrix"])
print(evaluation.statistical_tests)
\end{verbatim}

\section{Further reading}
\begin{thebibliography}{2}
\bibitem{mcc}
B. W. Matthews. Comparison of the predicted and observed secondary structure
of T4 phage lysozyme. \emph{Biochimica et Biophysica Acta}, 405(2), 1975.
\url{https://doi.org/10.1016/0005-2795(75)90109-9}.
\bibitem{reliability}
\software{} implementation and evaluation protocol:
\url{https://github.com/Jabualani/ML-Reliabilty}.
\end{thebibliography}
\end{document}
